\subsection{Retry: Repairing a Message Acceptance Condition}
\label{sec:retry}

Watch and Report observed execution. This project changes dispatch when
a request cannot yet be accepted. The question mark in the opening comment of
\filepath{examples/misc/fault-tolerant.abclr2} is worth keeping in mind.
This is a small experiment in reflective recovery, not a distributed
fault-tolerance service. Its receiver refuses ordinary requests until it
has received an initialization message. The sender's metaobject repairs
that condition and retransmits the refused request.
No machine crash or lost network packet is involved.

\subsubsection{Start with the Application}

Read the two application objects before their generator:

\wholelisting[firstline=130,lastline=147]{examples/misc/fault-tolerant.abclr2}

The server owns \code{amount}, initially \code{nil}. Its initialization
handler resets that value to zero, and \code{[:count-up n]} adds an amount
and prints the resulting count. The client has no application state:
\code{[:start server]} simply sends \code{[:count-up 10]} asynchronously.
Both explicitly select \code{FT-meta-gen}, so both participate in the
recovery protocol. The client's empty initialization handler is still
important: its metaobject must accept the client before it can run
\code{:start}.

In a fresh session at the repository root, the file's own entry sequence is:

\begin{codeblock}
\begin{lstlisting}
(inc "examples/misc/fault-tolerant" :use-source t)
[client <= [:initialize]]
[client <= [:start server]]
; Prints Count is 10.
\end{lstlisting}
\end{codeblock}

The server was deliberately not initialized by these commands. A successful
count therefore depends on recovery. Do not change \code{amount}'s initializer
to zero: that would conceal the problem the original example demonstrates.
For scratch-copy checks that leave the original directory untouched, run
\code{make -C tutorial retry-check}.

\subsubsection{Trace the Rejection and Retransmission}

Follow one request in six steps. This is a causal explanation, not a promise
about the order of unrelated console output.

\begin{enumerate}
\item The client's metaobject accepts \code{[:initialize]}, sets its private
acceptance flag, and executes the empty client handler.
\item The accepted \code{[:start server]} handler sends
\code{[:count-up 10]} to the server. Its sender is the client.
\item The server's metaobject dequeues the request. Its acceptance flag is
false, so it does not execute the arithmetic on \code{nil}.
\item It sends \code{[:reject Receiver Message Reply Sender]} to the
client's metaobject. Here \code{Receiver} is the server and
\code{Message} is the original \code{[:count-up 10]}.
\item The client's metaobject sends \code{[:initialize]} to the server,
then forwards the original message envelope to the server's metaobject.
\item The server processes initialization before the retransmitted request.
Initialization sets \code{amount} to zero; the accepted count request
then sets it to 10 and prints the count.
\end{enumerate}

The original recovery handler is short:

\wholelisting[firstline=47,lastline=51]{examples/misc/fault-tolerant.abclr2}

Notice the two levels of addressing. The initialization goes to the
application object \code{Receiver}. Retransmission goes directly to
\code{[meta Receiver]}, using the interpreter's \code{:message} envelope.
This preserves the old \code{Reply} and \code{Sender} fields rather than
creating a new application request with the recovering metaobject as sender.
In this trace the recovered request was a past-type send, so it has no
waiting application reply to deliver.

\subsubsection{Read the Acceptance State Machine}

The generator constructs the familiar queue, environment, script set,
evaluator, group manager, and \code{:dormant}/\code{:active} dispatch mode.
Its extra \code{initialized} flag belongs to the \emph{metaobject};
it is not the server's application variable \code{amount}.

\wholelisting[firstline=53,lastline=60]{examples/misc/fault-tolerant.abclr2}

The \code{:message} handler enqueues a three-field application envelope.
If dormant, the metaobject activates itself and schedules \code{:begin}.
That handler removes one envelope and applies the acceptance condition
before looking for an application script:

\wholelisting[firstline=66,lastline=80]{examples/misc/fault-tolerant.abclr2}

\code{[:initialize . \_]} matches a list beginning with the keyword
\code{:initialize}; the wildcard ignores its remaining tail. The flag is
set \emph{before} the application handler is found or executed. Thus the
comment saying that initialization has executed is stronger than the code:
the flag establishes that an initialization-shaped message was encountered.
It does not validate a successful initialization result.

For rejection, \code{Me} is the server's metaobject and \code{[den Me]}
recovers the application server it represents. \code{[meta Sender]}
selects the sender's recovery machinery. The rejected request has already
been removed from the queue; \code{:end} advances dispatch rather than
putting it back locally. The sender is responsible for retransmission.

For acceptance, lines 83--90 use \code{find-script} to obtain bindings and
a script body. The evaluator executes that body in an environment extending
the application's state. Its completion continuation schedules \code{:end}.
This continuation is an interpreter completion notice, not a client/server
acknowledgment. Lines 94--97 either process the next queued request or return
to dormant mode. The generator's final \code{:do-evlis} and \code{pairlis}
sequence builds the initial application environment, as in the earlier
inheritance and MCP generators; that construction is separate from the
application's \code{[:initialize]} protocol.

\subsubsection{What the Experiment Does and Does Not Guarantee}

A second client start prints \emph{Count is 20.}: the server is already
accepted, so no recovery reset is needed. An explicit server initialization
resets the count to zero; the next client start prints \emph{Count is 10.}
Six isolated checks verify these state transitions, the printed counts,
and an empty final server queue using the distributed example.
Reflective state queries inspect \code{amount}; they are test probes, not
new server API methods.

The example has no retry budget, timeout, persistent log, duplicate detection,
or exactly-once delivery mechanism. A sender without the expected
\code{:reject} handler cannot supply this recovery. Multiple rejections can
cause multiple initialization messages, each resetting server state.
The checked two-object sequence demonstrates one repaired precondition,
not correctness for every interleaving or failure.

\paragraph{Check your understanding, with answers.}
Why initialize the client even though its handler does nothing?
Its metaobject also gates application requests; the empty handler satisfies
the chosen protocol without adding application work.
Why preserve the original sender on retry?
Script matching can depend on the sender, and subsequent rejection must
reach the original sender's recovery machinery.
Would automatically resetting the server after every request be equivalent?
No: it would lose the accumulated count. Recovery belongs to the failed
acceptance path, not the normal successful path.
